##[简明数学史](智能简史之简明数学史.md)
##[现代数学的启发和计算机的发明](智能简史之现代数学的启发和计算机的发明.md)

###过去，人工智能从脑科学、心理学、经济学、量子力学中获取了能力。


在科学史上，无论是
桑代克的猫
巴普洛夫的狗
斯金纳的老鼠
……


###桑代克的饿猫迷笼实验
###斯金纳的老鼠
强化是塑造行为的有效而重要的条件。
塑造行为的过程，就是学习的过程。

实验1：每次按下按钮，则掉落食物。---正强化
结果：小白鼠自发学会了按按钮。
将行为与奖励不断重复、建立联系---强化学习的思路？

实验2【行为与惩罚】：将一只小白鼠放入一个有按钮的箱中。每次小白鼠不按下按钮，则箱子通电。---负强化
结果：小白鼠学会了按按钮。
但遗憾的是，一旦箱子不再通电，小白鼠按按钮的行为迅速消失。
惩罚可以迅速建立行为模式，但一旦惩罚消失，则行为模式也会迅速消失。

实验3【固定时间奖励】：由一直掉落食物，降低到每1分钟后按下按钮可掉落食物。
结果：小白鼠一开始不停按钮。过一段时间之后，小白鼠学会了间隔1分钟按一次按钮。
当掉落食物停止时，小白鼠的行为消失。

实验4【概率型奖励】：多次按下按钮，概率掉落食物。
结果：小白鼠学会了不停按钮。
当不再掉落食物时，小白鼠的学习行为消失速度非常慢。
为什么“赌博”会给予人类以依赖感/成瘾？

实验5【迷信的小白鼠】：概率型斯金纳箱。
结果：这些小白鼠有很多培养出了奇特的行为习惯，如撞箱子、转圈跳舞。
这是因为掉落食物前，小白鼠正好在进行这些行为，于是产生了“迷信”。
许多游戏中传出的谣言，比如“在中午抽奖容易得到大奖”之类，其原理与之相同。



###巴甫洛夫条件反射实验

学习：有机体在一定条件下形成刺激与反应的联系，从而获得新的经验。

对照：MachineLearning：AcomputerprogramissaidtolearnfromexperienceEwithrespecttosometaskTandsomeperformancemeasureP,ifitsperformanceonT,asmeasuredbyP,improveswithexperienceE.---TomMitchell,1998


###强化学习的多巴胺机制


###小结

过去，人工智能从脑科学、心理学、经济学、量子力学中获取了能力。
强化学习是一个类似多巴胺的机制，大多数时候很有效，但在某种情况下存在局限。
目前深度学习只是建立一个从输入到输出的映射过程，使错误最小化---只是数学上的最小化。
一旦AI在学习过程中发现某种路径最有利，就会陷入“局部最优解”而停滞不前。
表现得超乎寻常完美的神经网络其实并不真正理解其执行任务的意义。


